\documentclass[12pt]{article}
\usepackage{amsmath, amssymb, geometry, enumitem}
\geometry{margin=1in}
\setlength{\parindent}{0pt}
\setlength{\parskip}{0.2em}

\title{ECON 101: Macroeconomics \\ Solutions -- Quiz 3}
\author{Instructor: Muhebullah Karimzada}
\date{March 05,  2026}

\begin{document}
\maketitle

\section*{Part A: Multiple Choice Solutions}

\begin{enumerate}[label=\textbf{\arabic*.}]

\item \textbf{Correct answer: (C)} \; Greater technological innovation and productivity growth.\\
\textbf{Steps/Logic:}
\begin{itemize}
    \item Long-run growth in real GDP per person is driven mainly by \textbf{productivity} (output per worker/hour).
    \item Sustained productivity growth comes from \textbf{technological progress}, better organization, and accumulation of knowledge/human capital.
    \item (A) government consumption does not necessarily raise productive capacity; (B) reducing human capital hurts productivity; (D) lower participation reduces labour input and potential output.
\end{itemize}

\item \textbf{Correct answer: (B)} \; Real GDP per hour worked.\\
\textbf{Steps/Logic:}
\begin{itemize}
    \item Labour productivity measures how much output is produced per unit of labour input.
    \item In macro, labour input is commonly measured as \textbf{hours worked}.
    \item Therefore: \(\text{Labour productivity} = \dfrac{\text{Real GDP}}{\text{Hours worked}}\).
\end{itemize}

\item \textbf{Correct answer: (C)} \; Increase interest rates and crowd out private investment.\\
\textbf{Steps/Logic:}
\begin{itemize}
    \item A budget deficit means government needs to borrow more.
    \item In the loanable funds market, government borrowing increases \textbf{demand} for loanable funds.
    \item Higher demand (with supply unchanged) raises the \textbf{equilibrium real interest rate}.
    \item A higher interest rate reduces private borrowing for investment: this is \textbf{crowding out}.
\end{itemize}

\item \textbf{Correct answer: (B)} \; Channel savings efficiently into productive investment.\\
\textbf{Steps/Logic:}
\begin{itemize}
    \item Financial intermediaries (banks, credit unions, etc.) reduce transaction costs, screen/monitor borrowers, and allocate funds toward productive projects.
    \item This improves the efficiency of turning savings into investment, supporting capital accumulation and productivity growth.
    \item (A) intermediaries do not eliminate cycles; (C) and (D) are not primary roles.
\end{itemize}

\item \textbf{Correct answer: (B)} \; Investment equals national saving (closed economy).\\
\textbf{Steps/Logic:}
\begin{itemize}
    \item In a closed economy, there are no international capital flows.
    \item The resources available for investment come from \textbf{national saving}.
    \item Hence the accounting identity: \(\boxed{S = I}\) in a closed economy.
\end{itemize}

\end{enumerate}

\hrule
\vspace{1em}

\section*{Part B: Numerical Problem Solution}

\textit{Given:}
\[
Y = C + I + G
\]
\[
Y = 5{,}200,\quad C = 3{,}300,\quad G = 1{,}200
\]
(All values in billions of dollars.)

\subsection*{(a) Calculate investment \(I\) (3 marks)}

\textbf{Step 1: Start from the national income identity.}
\[
Y = C + I + G
\]

\textbf{Step 2: Rearrange to solve for \(I\).}
\[
I = Y - C - G
\]

\textbf{Step 3: Substitute the given values.}
\[
I = 5{,}200 - 3{,}300 - 1{,}200 = 700
\]

\textbf{Answer:}
\[
\boxed{I = 700 \text{ (billion dollars)}}
\]

\subsection*{(b) Compute national saving \(S\) (3 marks)}

\textbf{Key idea:} In a closed economy, \(\boxed{S = I}\).\\
(Reason: total saving must finance total investment when there is no foreign borrowing/lending.)

\textbf{Step 1: Use the result from part (a).}
\[
I = 700
\]

\textbf{Step 2: Apply the closed-economy identity.}
\[
S = I = 700
\]

\textbf{Answer:}
\[
\boxed{S = 700 \text{ (billion dollars)}}
\]

\textit{(Instructor note: If students know the decomposition \(S = (Y-T-C) + (T-G)\), they would need \(T\). Since taxes are not provided, the intended method is the closed-economy identity \(S=I\).)}

\subsection*{(c) Government spending rises by 200; taxes unchanged (4 marks)}

Government spending increases:
\[
\Delta G = +200
\]
Taxes are unchanged, so any change in the government budget balance comes from \(G\) only.

\textbf{Step 1: Effect on national saving.}
\begin{itemize}
    \item Government saving is \((T-G)\).
    \item If \(G\) rises and \(T\) is unchanged, then \((T-G)\) falls.
    \item Therefore national saving \(S\) falls.
\end{itemize}
In symbols (with unchanged \(T\)):
\[
\Delta (T-G) = -\Delta G = -200
\]
So:
\[
\boxed{\Delta S = -200}
\]
(i.e., national saving decreases by \$200 billion.)

\textbf{Step 2: Loanable funds market implication.}
\begin{itemize}
    \item The \textbf{supply} of loanable funds is national saving \(S\).
    \item A fall in \(S\) shifts the \textbf{supply curve of loanable funds left}.
\end{itemize}

\textbf{Step 3: Effect on the equilibrium real interest rate.}
\begin{itemize}
    \item With supply lower and demand unchanged, the equilibrium real interest rate rises.
\end{itemize}
\[
\boxed{\text{Equilibrium real interest rate increases.}}
\]

\textbf{Step 4: Crowding out (intuition, consistent with Ch. 6/7).}
\begin{itemize}
    \item Higher interest rates reduce private investment spending.
    \item This is the \textbf{crowding out} effect.
\end{itemize}

\textbf{Final conclusions for (c):}
\[
\boxed{S \text{ decreases by } 200,\quad r \text{ rises, and private } I \text{ is crowded out.}}
\]

\hrule
\vspace{1em}

\begin{center}
\textit{End of Solution Key}
\end{center}

\end{document}